\documentclass[11pt]{article}
\usepackage{setspace}
\usepackage{geometry}
\usepackage[hidelinks]{hyperref}
\geometry{margin=1in}
\setstretch{1.1}

\begin{document}

\begin{flushright}
\textbf{Zhongli Wang} \\
Phone: +49 160 9541 6015 \\
Email: \href{mailto:nkwangzhongli@gmail.com}{nkwangzhongli@gmail.com} \\
Website: \href{https://sites.google.com/view/zhongli-wang}{https://sites.google.com/view/zhongli-wang} \\
Bielefeld, Germany \\
\end{flushright}

\vspace{0.25cm}

\noindent
\textbf{Selection Committee} \\
Department of Computer Science \\
The Fu Foundation School of Engineering and Applied Science \\
Columbia University in the City of New York \\
New York, NY, USA \\

\vspace{0.25cm}

\noindent
\textbf{Application for Postdoctoral Researcher -- Enterprise AI}

\vspace{0.15cm}

\noindent Dear Members of the Selection Committee,\\[0.1cm]

I am writing to apply for the Postdoctoral Researcher position in the Department of Computer Science at Columbia University. I am close to completing my Ph.D.\ in Economics at Bielefeld University and Universit\`a Cattolica del Sacro Cuore within the Marie Sk{\l}odowska--Curie EPOC programme. My research sits at the intersection of agent-based modeling, machine learning, large-language-model agents, and computational social science. I am especially enthusiastic about building reliable, domain-grounded AI systems that can support knowledge work, research automation, decision-making, and enterprise-scale data analysis.

The Infosys Topaz Columbia University Enterprise AI Center is highly attractive to me because it combines rigorous academic research with practical system-building and real-world deployment. I am particularly drawn to the Center's focus on responsible, human-centered AI systems, enterprise automation, knowledge-grounded models, and AI adoption in organizational contexts. My background is not only in developing computational models, but also in translating complex research ideas into reproducible pipelines, prototypes, and applied tools. This makes the position an exciting opportunity for me to contribute to original research while working closely with Columbia faculty, Ph.D.\ students, and Infosys engineers.

A central strength I would bring to the Center is my hands-on experience with agentic AI workflows. In my recent co-authored project, Agentic Workflows for Economic Research, we develop a human-in-the-loop framework using LLM-based agents to support different stages of the research process, including literature search, conceptual synthesis, code generation, model implementation, and reproducible documentation. This project has given me direct experience with the practical challenges of building domain-specific agentic applications: task decomposition, prompt design, multi-agent coordination, human oversight, reliability, traceability, and evaluation of AI-assisted outputs. I see strong potential to extend this work toward enterprise use cases such as research intelligence, internal knowledge management, business analytics, compliance support, and AI-assisted decision workflows.

My doctoral dissertation further strengthens my ability to work on complex AI and data-driven systems. I develop computational models of digital market dynamics under consumer heterogeneity, network effects, congestion, and strategic firm investment. These models are implemented through large-scale simulation pipelines in Python and high-performance computing environments. In my third dissertation paper, I estimate an agent-based model using financial and operational data from AT\&T and T-Mobile, applying the Simulated Method of Moments and neural network-based surrogate modeling. This work has trained me to design scalable computational experiments, manage simulation outputs, evaluate model performance, and connect technical results to substantive business and policy questions.

Beyond my dissertation, my peer-reviewed article on Polytope Fraud Theory applies supervised machine learning to detect structured anomalies in financial statements. This project required feature engineering, model comparison, predictive evaluation, and interpretability-oriented analysis. It also sharpened my interest in trustworthy AI systems: models used in high-stakes organizational settings should not only perform well, but also provide transparent, auditable, and decision-relevant outputs. This perspective aligns closely with the Columbia--Infosys Center's emphasis on responsible AI, reliability, governance, and business impact.

Technically, I have strong experience with Python-based computational research, including simulation modeling, machine-learning workflows, data processing, visualization, and reproducible experimentation. I have also worked with Julia, R, C\#, SQL, Linux, Stata, MATLAB, EViews, Mathematica, and \LaTeX. My research has involved building agent-based models, running large-scale simulations, estimating structural and computational models, designing ML pipelines, and developing LLM-agent workflows. I am highly motivated to deepen this technical profile in AI systems, data management, knowledge-grounded models, and secure enterprise AI applications.

I would be particularly excited to contribute to the Center in three ways. First, I can help design and evaluate \textbf{domain-specific agentic systems} for enterprise knowledge work, with attention to reliability, reproducibility, and human oversight. Second, I can contribute to \textbf{knowledge-grounded AI and data-driven modeling} by combining structured domain knowledge with computational pipelines and empirical evaluation. Third, I can support \textbf{applied innovation and research translation} by building prototypes, preparing demonstrations, contributing to open-source releases, and communicating research outputs to both academic and industry audiences.

I also bring a collaborative and interdisciplinary working style. My doctoral training in the international EPOC network has required close collaboration across institutions, countries, and disciplines. I have presented my research at international conferences, organized and spoken at workshops on AI agents in science, and contributed to projects involving economists, data scientists, and computational modelers. I would be very glad to mentor students and interns, contribute to exploratory research directions, participate in proposal development, and represent the Center in academic and industry-facing venues.

Columbia University offers an exceptional environment for advancing responsible and practically useful AI, and I would be honored to contribute to the Infosys Topaz Columbia University Enterprise AI Center. I am open to relocation to the New York City area and would be excited to join a team that combines academic excellence, engineering ambition, and enterprise-scale impact.

Thank you so much for considering my application! I would be delighted to discuss how my background can contribute to the Center's research, prototypes, publications, and applied AI systems.

\vspace{0.35cm}

\noindent Sincerely,\\[0.95cm]
\noindent Zhongli Wang

\end{document}